\section{Policy Experiments}\label{sec:policy_experiments}

We use the estimated structural model to evaluate the behavioral, labor market, and economic consequences of introducing caregiving leave policies in Germany. The counterfactual policy experiments hold all model primitives---preferences, wage processes, health transitions, and care demand---fixed at their estimated values, and alter only the budget constraint and job offer process to reflect the new policy environment. This approach allows us to trace the general equilibrium of individual decisions from the moment a caregiving need arises through the remainder of the lifecycle.

\subsection{Policy Design}\label{subsec:policy_design}

We compare the existing German system of informal care cash benefits (\textit{Pflegegeld}) against two alternative caregiving leave policies, both of which include job retention guarantees. The three regimes are as follows.

\paragraph{Baseline: Pflegegeld.} The German long-term care insurance (\textit{Pflegeversicherung}) provides monthly cash benefits to informal caregivers: \euro{225} per month for light care needs and \euro{557.50} per month for intensive care needs (2010 values). These transfers are unconditional on labor supply---caregivers receive them regardless of whether they work full-time, part-time, or not at all. No job protection is provided; caregivers who reduce work or exit the labor force face normal job separation and job finding probabilities.

\paragraph{Normal caregiving leave (65\%, German Elterngeld-style).} This policy replaces the \textit{Pflegegeld} with an earnings-related leave benefit modeled on Germany's parental leave scheme. Eligible caregivers---those who provide informal care, are not retired, and had a job before the onset of caregiving---receive 65 percent of their previous net wage, bounded between \euro{300} and \euro{1,800} per month. The benefit is tax-free but subject to \textit{Progressionsvorbehalt} (\S 32b EStG), meaning it raises the marginal tax rate on other household income without itself being taxed. Crucially, the policy includes a job retention guarantee: while on caregiving leave, the agent faces zero job separation probability and is guaranteed a job offer upon returning to work. Experience accumulation continues as if the agent had remained in their pre-caregiving employment.

\paragraph{Full caregiving leave (100\%, Norwegian-style).} This more generous alternative provides 100 percent replacement of the caregiver's previous gross wage, capped at six times the Norwegian basic amount ($6G \approx$ \euro{56,700} annually in 2010). Unlike the normal leave, the full benefit is taxable income (though not subject to social security contributions). The same job retention guarantee and experience accumulation rules apply.

The key features shared by both leave policies---and absent from the baseline---are threefold: (i) the benefit is tied to prior earnings, creating stronger incentives for higher-wage agents; (ii) the job retention guarantee eliminates the risk of permanent labor market exit during caregiving; and (iii) experience continues to accumulate, preserving human capital and future pension entitlements. These features jointly alter the intertemporal trade-off agents face when a parent requires care.

Table~\ref{tab:policy_behavioral_changes} reports the complete set of behavioral outcomes across the three regimes, organized into thirteen panels covering caregiving arrangements, labor supply composition, economic outcomes, and retirement. We discuss the key findings in turn.


\subsection{Effects on Caregiving Arrangements}\label{subsec:cg_arrangements}

\paragraph{Aggregate effects.} Both leave policies induce a substantial reallocation of care provision from formal to informal arrangements. Panel~B of Table~\ref{tab:policy_behavioral_changes} reports caregiving shares conditional on the parent having a positive care need. Under the baseline, 42.6 percent of agents with a care-dependent parent choose informal care, while 33.4 percent choose formal care. The normal leave policy raises the informal care share by 8.0 percent (to 46.1 percent), while the full leave increases it by 9.9 percent (to 46.9 percent). The increase in informal care is almost exactly mirrored by a decrease in formal care, which falls by 10.2 percent under normal leave and 12.6 percent under full leave. In percentage-point terms, both flows amount to approximately 3.4 and 4.2 percentage points, respectively---a near-perfect one-for-one substitution between formal and informal arrangements.

The increase in informal care is distributed across both care intensities. Light informal care rises by 7.1 percent (normal leave) and 8.9 percent (full leave), while intensive informal care---which includes combination care, where the caregiver supplements informal provision with formal home care services---rises by 9.2 and 11.1 percent, respectively. The somewhat larger response of intensive care suggests that the leave policies particularly encourage agents to take on more demanding caregiving roles that they would have delegated to formal providers under the baseline.

The age pattern of the caregiving response is strongly concentrated in the 50--59 age range. Among agents aged 55--59, informal care rises by 13.4 percent (normal leave) and 15.7 percent (full leave), while formal care falls by 16.6 and 19.4 percent---the largest responses of any age group. This age range represents the intersection of peak care demand (parents are typically in their late 70s to 80s) and peak earning potential, making the opportunity cost of informal care highest and the leave benefit most valuable in absolute terms. At ages 40--49, where care demand is less prevalent and wages lower, the response is more muted (4.5 and 6.0 percent). At 60--67, the response is again smaller (4.7 and 6.5 percent), as agents in this range are approaching retirement and face a different set of margins.

\paragraph{Education gradient in the caregiving response.} The behavioral response to the leave policies is strongly heterogeneous across education groups, consistent with the opportunity cost mechanism embedded in the model. Panels~C and~D of Table~\ref{tab:policy_behavioral_changes} decompose the caregiving results by education level.

Under normal leave, high-education agents increase informal care by 11.4 percent, nearly twice the 6.6 percent increase among low-education agents. Under full leave, the disparity widens further: high-education agents increase informal care by 15.7 percent compared to 7.3 percent for their low-education counterparts. The formal care response mirrors this gradient, with high-education agents reducing formal care use by 14.5 percent (normal) and 20.0 percent (full), versus 8.4 and 9.4 percent for low-education agents.

The education gradient is sharpest among agents aged 55--59, where the interaction of high wages and intense care demand amplifies the policy response. High-education agents in this age range increase informal care by 20.3 percent under normal leave and 25.7 percent under full leave---more than double the response of low-education agents in the same age group (10.6 and 11.6 percent). At the same time, high-education agents reduce formal care by 22.2 and 28.1 percent, compared to 13.8 and 15.1 percent for low-education agents.

The economic mechanism driving this differential is straightforward. High-education agents earn higher wages, implying a higher opportunity cost of forgoing market work to provide informal care. Under the baseline \textit{Pflegegeld}---a flat transfer independent of earnings---these agents rationally opt for formal care. The leave policies, by contrast, replace a fraction of prior earnings, providing a larger absolute benefit to higher earners. The 65 percent replacement rate of the normal leave and the 100 percent rate of the full leave both deliver a greater absolute subsidy to high-education agents, reducing their effective opportunity cost of informal care by more. The result is a disproportionate behavioral shift toward informal care among high earners.

The composition of the caregiving response also differs by education. Among low-education agents, who already provide intensive care at high rates (21.6 percent versus 12.6 percent for high-education agents), the policy primarily intensifies existing patterns---the intensive care response is 7.8 percent. Among high-education agents, who have low baseline intensive care provision, the policy induces a much larger relative shift into intensive care (14.8 percent under normal leave, 20.6 percent under full leave). The leave benefit thus induces high-education agents to take on caregiving roles they would not have chosen in the absence of the policy.

\paragraph{Sensitivity to policy generosity.} The differential response between normal and full leave is itself informative about the elasticity of caregiving choices with respect to the replacement rate. For low-education agents, the increase from 65 to 100 percent replacement adds only 0.7 percentage points to the informal care response (from 6.6 to 7.3 percent). For high-education agents, the same increase in generosity adds 4.3 percentage points (from 11.4 to 15.7 percent). This asymmetry reflects the fact that the gap between 65 and 100 percent replacement translates into a substantially larger absolute benefit differential for high earners. Low-education agents, whose prior wages are lower, are already near the policy's lower bound under normal leave and gain relatively little from the move to full replacement.

This pattern has implications for optimal policy design: the marginal fiscal cost of increasing generosity from 65 to 100 percent is borne disproportionately by subsidizing high earners' informal care provision, while low-education agents---who already bear a heavier caregiving burden in the baseline---see comparatively smaller behavioral changes.


\subsection{Labor Market Responses}\label{subsec:labor_market}

\paragraph{Population-level effects are modest.} The aggregate labor market consequences of the leave policies are small. Panel~A of Table~\ref{tab:policy_behavioral_changes} shows that the overall employment rate among agents aged 30--67 changes by less than 0.1 percent under either policy. Full-time and part-time employment shares are similarly stable. The only notable population-level effect is a 1.7 percent reduction in the retirement share among agents aged 60--67 under normal leave (1.8 percent under full leave), corresponding to a decline from 40.9 to 40.2 percent. These small aggregate effects reflect the fact that the policies directly target only the subset of agents who face parental care needs, and the behavioral changes within this group are partially offset across the lifecycle.

\paragraph{During active caregiving: large employment reductions.} The labor market effects are far more dramatic when we condition on agents who are currently providing informal care. Panel~E shows that the employment rate among active caregivers falls from 57.9 percent in the baseline to 49.4 percent under normal leave and 48.6 percent under full leave---declines of 14.8 and 16.1 percent, respectively. Full-time work among caregivers falls by 16.6 percent (normal) and 22.5 percent (full), while the share classified as unemployed---which in this context captures leave take-up---surges by 38.5 and 42.5 percent.

These responses are economically large. Under the full leave, the share of caregivers who stop working entirely rises from 28.8 to 41.0 percent, indicating substantial take-up of the policy. The ``unemployment'' category in the model captures agents who cease market work to devote time to caregiving while receiving the leave benefit; the sharp increase in this category is the behavioral manifestation of the leave incentive.

The within-caregiving employment response is strongly concentrated in the 50--59 age range, where employment falls by 21.1 percent (normal) and 22.6 percent (full). Among agents aged 55--59, the effect reaches 23.4 and 25.3 percent. At ages 60--67, the response is more muted but still meaningful at 5.8 and 8.0 percent. Notably, the retirement share among caregivers aged 60--67 drops sharply, from 37.9 to 31.8 percent under normal leave ($-16.3$ percent) and 30.7 percent under full leave ($-19.2$ percent). The leave policy thus keeps older caregivers attached to the labor market: rather than retiring early to provide care, they take temporary leave while maintaining their employment relationship.

\paragraph{The transition into caregiving.} Panel~K examines the labor market at the onset of caregiving, when agents first begin providing informal care. The employment rate at the first caregiving year falls by 4.2 percent (normal) and 4.3 percent (full) relative to the baseline. The full-time share declines by 6.2 percent under normal leave and 11.1 percent under full leave, while part-time employment actually rises slightly under full leave (+1.1 percent). This pattern suggests that agents substitute from full-time to part-time work rather than exiting the labor force entirely, taking advantage of the gap payment available to full-time-to-part-time transitions under the full leave.

The asymmetry between the two policies at the first caregiving year is informative. The full-time decline is nearly twice as large under full leave ($-11.1$ versus $-6.2$ percent), indicating that agents are more willing to reduce hours when the replacement rate fully compensates their earnings loss. The overall employment rate decline, however, is nearly identical across policies ($-4.2$ versus $-4.3$ percent), implying that the additional generosity of the full leave primarily shifts agents along the intensive margin (full-time to part-time) rather than the extensive margin (working to not working).

\paragraph{Lifecycle employment: near-zero net effects.} Despite the substantial labor market disruption during active caregiving, the lifetime employment of ever-caregivers is virtually unchanged by either policy. Panels~F and~L of Table~\ref{tab:policy_behavioral_changes} report the employment composition of ever-caregivers across all periods of their lifecycle. The overall employment rate changes by less than 0.2 percent, whether we include all periods (Panel~F) or exclude the first caregiving year (Panel~L).

This near-zero lifecycle effect conceals a precisely offsetting pattern across the age distribution. At ages 30--49, before most caregiving episodes begin, employment effects are negligibly positive (+0.1 to +0.3 percent). During the peak caregiving window at ages 50--59, employment falls by 1.8 to 2.4 percent---a direct consequence of the leave take-up documented above. At ages 60--67, however, employment rises by 4.2 percent (normal) and 4.1 percent (full), driven by a 3.7 and 4.2 percent reduction in retirement at these ages.

The interpretation is clear: the leave policies redistribute the timing of labor supply across the lifecycle without altering its total quantity. Agents reduce work during the caregiving episode but compensate by working longer at older ages, enabled by the job retention guarantee that preserves their access to employment. The result is a temporal reallocation of work, not a net reduction.


\subsection{The Role of Caregiving Duration}\label{subsec:duration}

The lifecycle labor supply effects documented above mask substantial heterogeneity by the duration of the caregiving episode. Panel~M of Table~\ref{tab:policy_behavioral_changes} classifies ever-caregivers by the total number of years they provide informal care---one year, two to three years, or four or more years---and reports labor supply outcomes for each group separately. The resulting duration gradient is among the most striking findings of the policy analysis.

\paragraph{Short-duration caregivers are barely affected.} Agents who provide informal care for a single year experience negligible changes in lifetime employment ($-0.3$ percent) and a slight increase in retirement at ages 63--67 (+0.4 percent under normal leave). The caregiving episode is too brief to generate meaningful behavioral responses: by the time the leave benefit takes effect, the caregiving need has subsided and the agent returns to their prior labor market state. The policy is essentially irrelevant for this group.

\paragraph{Moderate-duration caregivers show small effects.} Agents with two to three years of caregiving exhibit near-zero lifetime employment effects and a small reduction in retirement ($-1.5$ percent under normal leave, $-1.8$ percent under full leave at ages 63--67). The leave policy begins to matter for this group, keeping them attached to the labor force through the caregiving spell and slightly delaying retirement, but the effects remain modest.

\paragraph{Long-duration caregivers benefit substantially.} The most consequential effects emerge among agents who provide informal care for four or more years. For this group, the leave policy \textit{increases} lifetime employment by 1.0 percent (normal) and 1.1 percent (full)---the opposite sign of the within-caregiving effect. The mechanism is a dramatic reduction in early retirement: the share of long-duration caregivers who are retired at ages 63--67 falls by 9.7 percent under normal leave and 10.6 percent under full leave, from a baseline of 72.3 percent.

The age-specific employment effects for long-duration caregivers are especially revealing. At ages 60--67, the employment rate increases by 15.5 percent under both policies---from 32.6 to 37.7 percent. At ages 50--59, the peak caregiving window, employment falls by 3.2 to 3.5 percent as agents take up the leave benefit. The combination of temporary labor force withdrawal during caregiving and sustained labor force attachment afterward produces the positive net lifetime effect.

The economic interpretation is as follows. Under the baseline, agents who face multi-year caregiving responsibilities---arising from parents with persistent health deterioration---confront a stark choice: continue working while outsourcing care to formal providers, or exit the labor force to provide care themselves. Without job protection, those who choose informal care face elevated job separation risk, making re-entry after caregiving difficult or impossible. Many respond by retiring early, permanently forgoing future earnings and pension accumulation. The leave policy transforms this decision by offering temporary income replacement with a guaranteed path back to employment. Agents who would have retired permanently instead take leave, provide care for their parent, and return to work---preserving both their income trajectory and their pension entitlements.


\subsection{Experience Accumulation and Pension Outcomes}\label{subsec:pensions}

The job retention and experience accumulation provisions of the leave policies have meaningful consequences for pension entitlements and retirement timing. Panel~J of Table~\ref{tab:policy_behavioral_changes} reports experience years and retirement age under each regime.

Under the normal leave, ever-caregivers accumulate 3.0 percent more experience years by the time they retire (3.3 percent under full leave), corresponding to an increase of approximately 0.74 years (0.82 years under full leave) of credited work experience. For the full population of retirees---including those who never provided care---the experience gain is 1.3 percent (1.4 percent under full leave), or roughly one-third of a year. These gains arise because the leave policies credit experience during the caregiving episode as if the agent had continued in their pre-caregiving job, whereas under the baseline, agents who exit the labor force to provide care accumulate no experience.

The experience gains translate directly into higher pension incomes. Panel~I shows that average gross pension income rises by 1.3 percent under normal leave (\euro{134} per year) and 1.4 percent under full leave (\euro{148} per year). While modest in percentage terms, these gains are permanent---they accrue annually for the remainder of the retiree's life---and concentrated among ever-caregivers, for whom the proportional pension gains are larger.

The retirement age itself shifts upward, consistent with the delayed retirement documented in Section~\ref{subsec:duration}. Average retirement age for ever-caregivers rises by approximately 0.13 years (roughly 1.5 months) under normal leave and 0.14 years under full leave. This delay reflects both the direct effect of the job retention guarantee---agents who would have retired to provide care instead take leave---and the indirect effect of higher accumulated experience: agents with more experience have higher continuation values from working, further incentivizing delayed retirement.

Panel~H reports pre-retirement economic outcomes for agents under age 63. Average gross labor income declines by approximately 0.2 percent under both policies, reflecting the temporary work reduction during caregiving episodes. Average wealth declines marginally (0.1--0.3 percent), consistent with lower savings during the caregiving window. Average consumption, however, is virtually unchanged under normal leave and actually rises by 0.2 percent under full leave, suggesting that the leave benefit successfully smooths consumption through the caregiving episode. The full leave's consumption-preserving effect is consistent with its more generous replacement rate providing a closer approximation to permanent income during the temporary disruption.


\subsection{The Structure of Transfers}\label{subsec:benefits}

The structure of benefits received by caregivers differs markedly across regimes and reveals the incentive mechanism underlying the behavioral responses documented above. Panel~G of Table~\ref{tab:policy_behavioral_changes} reports average benefits by labor market state during caregiving.

Under the baseline, all informal caregivers receive the flat \textit{Pflegegeld} regardless of their employment status, averaging approximately \euro{4,100}--\euro{4,500} per year across labor states and age groups. Both leave policies replace the \textit{Pflegegeld} with the earnings-related benefit, creating sharp differences in transfers by labor market state.

Caregivers who continue working full-time receive no leave benefit under either policy---the benefit compensates only for earnings losses---and thus forfeit the \euro{4,061} annual \textit{Pflegegeld} they would have received under the baseline. Part-time caregivers receive \euro{444} per year under normal leave (far less than the \euro{4,235} baseline) but \euro{7,044} under full leave (substantially more). The largest benefits accrue to caregivers who stop working entirely to provide care: they receive \euro{7,900} under normal leave and \euro{16,855} under full leave, compared to \euro{5,186} under the baseline.

This benefit structure reveals a fundamental reorientation of the transfer incentive. The \textit{Pflegegeld} provides a flat subsidy independent of labor supply, creating no marginal incentive to adjust work behavior. The leave policies, by contrast, concentrate transfers on agents who reduce their labor supply to provide care, generating a strong incentive to move from employment to leave. The large employment reductions during active caregiving documented in Section~\ref{subsec:labor_market} are the behavioral consequence of this incentive structure: the leave policies effectively subsidize the act of withdrawing from the labor market to provide care, whereas the \textit{Pflegegeld} does not.


\subsection{Discussion}\label{subsec:discussion}

The policy experiments reveal that caregiving leave policies operate through two distinct channels with different time horizons and distributional properties.

\paragraph{The short-run channel: substitution from formal to informal care.} In the short run, the leave policies reduce the effective price of informal care provision, inducing a reallocation from formal to informal arrangements. This reallocation is approximately one-for-one: each percentage point of additional informal care is matched by a percentage point reduction in formal care. The substitution is concentrated in the 50--59 age range, where the intersection of high care demand and high wages creates the largest behavioral margin. The short-run employment effects are large and negative---active caregivers reduce employment by 15--16 percent---with the freed-up time redirected from market work to care provision.

\paragraph{The long-run channel: job retention and delayed retirement.} Over the lifecycle, the negative employment effects during caregiving are almost exactly offset by positive effects at older ages. The job retention guarantee prevents permanent labor force exit, enabling agents to return to work after the caregiving episode concludes. This mechanism is most consequential for long-duration caregivers (four or more years), who under the baseline would retire permanently to provide extended care. With the leave policy, these agents instead take temporary leave and return, generating a 10 percent reduction in early retirement and a 15 percent increase in employment at ages 60--67. The experience accumulation provision further supports this channel by preserving human capital through the caregiving spell, producing 3 percent more experience at retirement and commensurately higher pension entitlements.

\paragraph{Distributional implications.} The education gradient in the caregiving response highlights an important distributional dimension of earnings-related leave policies. High-education agents, who face higher opportunity costs of informal care, respond nearly twice as strongly as low-education agents. Under the full leave, high-education agents increase informal care by 15.7 percent and reduce formal care by 20.0 percent, compared to 7.3 and 9.4 percent for low-education agents. Since the benefit is proportional to prior earnings, the policy transfers are also larger in absolute terms for high earners. The primary behavioral effect of the policy thus operates through the group with the highest wages and lowest baseline informal care provision, while the group that already bears the heavier caregiving burden---low-education agents---experiences comparatively smaller behavioral changes.

\paragraph{Normal versus full leave.} The comparison between the 65 percent and 100 percent replacement rates provides insight into the elasticity of caregiving decisions with respect to policy generosity. The additional 35 percentage points of replacement add approximately 2 percentage points to the overall informal care response (from 8.0 to 9.9 percent) and a comparable amount to the formal care reduction. The marginal effect of increased generosity is concentrated among high-education agents (4.3 additional percentage points of informal care) and is negligible for low-education agents (0.7 percentage points). The full leave produces modestly larger effects on labor supply during caregiving---a 22.5 percent decline in full-time work among caregivers versus 16.6 percent under normal leave---but the lifecycle employment effects are nearly identical across the two policies. This suggests that the job retention guarantee, common to both policies, is the primary driver of long-run outcomes, while the replacement rate principally governs the short-run intensity of leave take-up and the composition of care.

\paragraph{Implications for policy design.} These findings point to several design lessons. First, the job retention guarantee is the key policy lever for long-run labor market outcomes. Without it, caregiving leave would generate substantial employment disruptions during the caregiving episode with no offsetting gains at older ages; with it, the leave policy preserves labor market attachment and converts a permanent exit into a temporary absence.

Second, the leave policies create both winners and losers among caregivers relative to the existing \textit{Pflegegeld}. The earnings-related benefit concentrates transfers on agents who reduce their labor supply; caregivers who continue working full-time---26.3 percent of informal caregivers in the baseline---lose their \textit{Pflegegeld} and receive no leave benefit. Under the normal leave, only the 28.8 percent of caregivers who stop working entirely receive more than under the baseline, while the average caregiver receives \euro{1,128} less per year in direct transfers (\euro{3,341} versus \euro{4,469}). Under the full leave, the average transfer rises above the \textit{Pflegegeld} (\euro{9,091}), but full-time working and retired caregivers still receive nothing. The normal leave's advantage over the \textit{Pflegegeld} is therefore not in transfer generosity but entirely in the job retention and experience accumulation channels.

Third, the earnings-related structure compounds the education gradient documented above: high-education agents, who have higher prior wages, receive larger absolute leave benefits, while low-education agents---who already bear a disproportionate caregiving burden in the baseline---receive smaller transfers and exhibit smaller behavioral responses. The leave policy's primary fiscal expenditure thus flows to high-education agents transitioning from formal to informal care, rather than to low-education agents who already provide informal care without the incentive. A policy that combines a flat, unconditional caregiving benefit with the job retention guarantee would preserve the long-run labor market benefits documented here while providing more equitable support across the earnings distribution.

Fourth, the distinction between 65 and 100 percent replacement is primarily relevant for high-education agents and for the short-run composition of labor supply. The long-run employment and pension effects are largely insensitive to the replacement rate, driven instead by the job retention mechanism that both policies share.
